% Generated draft: requires mathematical and source review.
\documentclass[11pt]{article}
\usepackage[margin=25mm]{geometry}
\usepackage{fontspec}
\setmainfont{Latin Modern Roman}
\usepackage{amsmath,amssymb,mathtools,unicode-math}
\setmathfont{Latin Modern Math}
\usepackage{enumitem,needspace,xurl,hyperref}
\setlength{\parindent}{0pt}
\setlength{\parskip}{5pt}
\setlength{\emergencystretch}{3em}
\newcommand{\N}{\mathbb{N}}
\newcommand{\Z}{\mathbb{Z}}
\newcommand{\Q}{\mathbb{Q}}
\newcommand{\R}{\mathbb{R}}
\newcommand{\C}{\mathbb{C}}
\newcommand{\F}{\mathbb{F}}
\newcommand{\A}{\mathbb{A}}
\newcommand{\PP}{\mathbb{P}}
\newcommand{\E}{\mathbb{E}}
\newcommand{\eps}{\varepsilon}
\newcommand{\dd}{\,\mathrm{d}}
\begin{document}
\section*{UnsolvedMath ID: 30006661}
\textbf{Source problem code:} catalogue entry
\subsection{Definitions and mathematical statement}
Let \(X\) and \(Y\) be finite two-dimensional CW complexes, subject to the supplied hypothesis that their fundamental groups are finite. Here finite means that there are finitely many cells. A homotopy equivalence \(X\simeq Y\) consists of continuous maps \(f:X\to Y\) and \(g:Y\to X\) such that \(g\circ f\simeq\operatorname{id}_X\) and \(f\circ g\simeq\operatorname{id}_Y\).

Write \(M(X)\) and \(M(Y)\) for the smooth four-dimensional doubles referred to in the supplied record. For context, the smooth double of a compact smooth four-manifold \(W\) with boundary is obtained by gluing two copies of \(W\) along their boundaries by the identity, using collars to give the result a smooth structure. The supplied record does not specify the thickenings, choices, or precise construction intended by \(M(X)\) and \(M(Y)\); the general description of doubling therefore does not by itself define these assignments.

\textbf{Question (permanent catalogue entry).} For every pair \(X,Y\) satisfying the supplied hypotheses, does
\[
X\simeq Y\quad\Longrightarrow\quad M(X)\cong_{\mathrm{diff}}M(Y)
\]
hold, where \(\cong_{\mathrm{diff}}\) means that there exists a smooth bijection with smooth inverse?

\textbf{Unresolved qualifications.} This is the question in the supplied reconstructed statement, conditional on clarification of the double construction and its permitted choices. The record does not state whether the conclusion is to hold for every choice of thickening or for particular prescribed choices. It also does not explicitly require connectedness, specify basepoints for the fundamental groups, or specify an orientation requirement for the diffeomorphism. These conventions require confirmation from the source. The supplied original extract contains the label ``Question 2.2'' followed by a discussion of the quadratic bias invariant and Tate cohomology, but does not contain the displayed question above.
\subsection{Short English statement}
The supplied question asks whether homotopy-equivalent finite two-dimensional CW complexes with finite fundamental group have diffeomorphic smooth four-dimensional doubles. The intended construction of the doubles and the quantification over construction choices require clarification. This is a statement-only draft requiring human mathematical review.
\subsection{Sources}
\begin{itemize}
\item Supplied local UnsolvedMath registry record: permanent ID \(30006661\), title \emph{Diffeomorphic Doubles of Homotopy-Equivalent Two-Complexes}, source code \texttt{catalogue entry}, in the Oberwolfach Reports Open Problems collection. The record gives proposed year 2026 and names no proposer.
\item \emph{Cohomology of Finite Groups: Interactions and Applications} (2026), Oberwolfach Reports, pp. 541--610, DOI 10.4171/owr/2026/9. Supplied report link: \url{https://doi.org/10.4171/owr/2026/9}. The supplied original extract includes page marker 605, the label ``Question 2.2'', and surrounding discussion of the quadratic bias invariant and Tate cohomology; its mathematical notation is damaged and the question text is missing.
\item Hambleton--Nicholson, \emph{Four-manifolds, two-complexes and the quadratic bias invariant}. Supplied literature link: \url{https://doi.org/10.1016/j.aim.2025.110728}.
\item Supplied literature assessment, checked on \texttt{2026-08-22}: the record labels the question open and states that the quadratic bias invariant provides new obstructions and examples but no general diffeomorphism theorem. This is supplied dated curation derived from the Oberwolfach report, not an independently verified assessment; time-sensitive status and cited claims require independent verification.
\item The registry labels the statement \texttt{corrected\_verified} and reports that it was reconstructed and checked against the official Oberwolfach report. That verification claim is supplied provenance, not verification performed for this draft. The report and paper above are supplied citations; neither was independently consulted here.
\end{itemize}
\end{document}
